%%%%%%%%%%%%%%%%%%%%%%%%%%%%%%%%%%%%%%%%%%%%%%
%
%		Thesis Settings
%
%		UTV Template
%		2024
%
%%%%%%%%%%%%%%%%%%%%%%%%%%%%%%%%%%%%%%%%%%%%%%

\usepackage[T1]{fontenc}
\usepackage[utf8]{inputenc}
\usepackage[english]{babel}

\usepackage[left=1.5cm,right=1.5cm,top=1.5cm,bottom=2cm]{geometry}
\usepackage{fourier} % Utopia font-typesetting including mathematical formula compatible with newer TeX-Distributions (>2010)
% \usepackage{utopia} % On older systems -> use this package instead of fourier in combination with mathdesign for better looking results
% \usepackage[adobe-utopia]{mathdesign}

\usepackage{setspace}

\makeatletter
\setlength{\@fptop}{0pt}  % For aligning all floating figures/tables etc. to the top margin
\makeatother

\usepackage{graphicx,xcolor}
\graphicspath{{images/}}

\usepackage{subfig}
\usepackage{booktabs}
\usepackage{lipsum}
\usepackage{microtype}
\usepackage{url}
\usepackage[final]{pdfpages}
\usepackage[shortlabels]{enumitem}
\usepackage{color}

\usepackage{tikz}
\usetikzlibrary{positioning}
\usetikzlibrary{calc}
\usetikzlibrary{arrows.meta}

%%%%% PAGE HEADER %%%%
\usepackage{fancyhdr}
% 1) Define how LaTeX stores chapter and section info in the marks:
\renewcommand{\chaptermark}[1]{%
  \markboth{\chaptername\ \thechapter\quad #1}{}%
}
\renewcommand{\sectionmark}[1]{%
  \markright{\thesection\quad #1}%
}

% 2) Activate fancyhdr:
\pagestyle{fancy}
\fancyhf{} % clear all headers and footers

% 3) Rule widths:
\renewcommand{\headrulewidth}{0.4pt}
\renewcommand{\footrulewidth}{0pt}

% 4) Define headers:
% EVEN (LEFT) PAGES
% Header: Chapter info on the left, nothing on the right
\fancyhead[LE]{\bfseries \nouppercase{\leftmark}}
\fancyhead[RE]{}
% Footer: Page number on the left, nothing on the right
\fancyfoot[LE]{\thepage}
\fancyfoot[RE]{}

% ODD (RIGHT) PAGES
% Header: Nothing on the left, section info on the right
\fancyhead[LO]{}
\fancyhead[RO]{\bfseries \nouppercase{\rightmark}}
% Footer: Nothing on the left, page number on the right
\fancyfoot[LO]{}
\fancyfoot[RO]{\thepage}

% 5) Redefine plain style for chapter start pages:
\fancypagestyle{plain}{%
  \fancyhf{}%
  \renewcommand{\headrulewidth}{0pt}%
  \renewcommand{\footrulewidth}{0pt}%
  % Place page number at outer foot (left for even pages, right for odd pages)
  \fancyfoot[LE,RO]{\thepage}%
}

\usepackage{listings}
% \lstset{language=[LaTeX]Tex,tabsize=4, basicstyle=\scriptsize\ttfamily, showstringspaces=false, numbers=left, numberstyle=\tiny, numbersep=10pt, breaklines=true, breakautoindent=true, breakindent=10pt}
\lstdefinestyle{oneside}{
    basicstyle=\ttfamily\scriptsize,
    breakatwhitespace=false,
    tabsize=2,
    showtabs=false,
    showstringspaces=false,
    showspaces=false,
    numbers=left,
    numberstyle=\tiny,
    numbersep=5pt,
    keepspaces=true,
    captionpos=b,
    frame=lines
}
\lstset{style=oneside}

\usepackage{hyperref}
\hypersetup{pdfborder={0 0 0},
	colorlinks=true,
	linkcolor=black,
	citecolor=black,
	urlcolor=black}
\urlstyle{same}

\makeatletter
\def\cleardoublepage{\clearpage\if@twoside \ifodd\c@page\else
    \hbox{}
    \thispagestyle{empty}
    \newpage
    \if@twocolumn\hbox{}\newpage\fi\fi\fi}
\makeatother \clearpage{\pagestyle{plain}\cleardoublepage}

%%%%% CHAPTER HEADER %%%%
\usepackage[explicit]{titlesec}
\newcommand*\chapterlabel{}
%\renewcommand{\thechapter}{\Roman{chapter}}
\titleformat{\chapter}[display]  % Type (section,chapter,etc...) to vary,  shape (e.g. display-type)
	{\normalfont\bfseries\huge} % Format of the chapter
	{\gdef\chapterlabel{\thechapter\ }} % The label
 	{0pt} % Separation between label and chapter-title
 	{%
    \begingroup
    \microtypesetup{expansion=false, protrusion=false}%
    \hyphenpenalty=10000
    \exhyphenpenalty=10000
    {\parskip=0pt
      \begin{tikzpicture}[remember picture,overlay]
        \node[yshift=-8cm] at (current page.north west)
          {%
            \begin{tikzpicture}[remember picture, overlay]
              \draw[fill=black] (0,0) rectangle(25mm,10mm);
              \node[anchor=north east,yshift=-7.23cm,xshift=24mm,minimum height=30mm,inner sep=0mm] at (current page.north west)
              {\parbox[top][30mm][t]{15mm}{\raggedleft $\phantom{\textrm{l}}$\color{white}\chapterlabel}}; % The black l is just to get better base-line alignment
              \node[anchor=north west,yshift=-7.23cm,xshift=27mm,text width=\textwidth,minimum height=30mm,inner sep=0mm] at (current page.north west)
              {\parbox[top][30mm][t]{\textwidth}{\color{black}#1}};
            \end{tikzpicture}
          };
      \end{tikzpicture}
    }
    \gdef\chapterlabel{}
    \endgroup
  } % Code before the title body

\titlespacing*{\chapter}{0pt}{110pt}{40pt}
\titlespacing*{\section}{0pt}{12pt}{6pt}
\titlespacing*{\subsection}{0pt}{12pt}{6pt}
\titlespacing*{\subsubsection}{0pt}{12pt}{6pt}

\usepackage{array}
\usepackage{amsmath}
\usepackage{amssymb}
\usepackage{amsthm}
\usepackage{amsfonts}
\usepackage{mathtools}

\usepackage{algorithm}
\usepackage{algpseudocode}
\usepackage{float}
\floatstyle{ruled}
\restylefloat{algorithm}
\algrenewcommand\algorithmicrequire{\textbf{Input:}}
\algrenewcommand\algorithmicensure{\textbf{Output:}}

% Fix the problem with delimiter size caused by fourier and amsmath packages.
\makeatletter
\def\resetMathstrut@{%
  \setbox\z@\hbox{%
    \mathchardef\@tempa\mathcode`\(\relax
      \def\@tempb##1"##2##3{\the\textfont"##3\char"}%
      \expandafter\@tempb\meaning\@tempa \relax
  }%
  \ht\Mathstrutbox@1.2\ht\z@ \dp\Mathstrutbox@1.2\dp\z@
}
\makeatother

\usepackage{subfiles} % Compile sub-main files

\usepackage[capitalise,sort]{cleveref}

\newtheorem{theorem}{Theorem}[chapter]
\newtheorem{assumption}[theorem]{Assumption}
\newtheorem{definition}[theorem]{Definition}
\newtheorem{problem}[theorem]{Problem}
\newtheorem{lemma}[theorem]{Lemma}
\newtheorem{claim}[theorem]{Claim}
\newtheorem{remark}[theorem]{Remark}
\newtheorem{corollary}[theorem]{Corollary}
\newtheorem{proposition}[theorem]{Proposition}
\newtheorem{example}[theorem]{Example}

\crefname{theorem}{Theorem}{Theorems}
\crefname{assumption}{Assumption}{Assumptions}
\crefname{definition}{Definition}{Definitions}
\crefname{problem}{Problem}{Problems}
\crefname{lemma}{Lemma}{Lemmas}
\crefname{claim}{Claim}{Claims}
\crefname{remark}{Remark}{Remarks}
\crefname{corollary}{Corollary}{Corollaries}
\crefname{proposition}{Proposition}{Propositions}
\crefname{example}{Example}{Examples}
\crefname{algorithm}{Algorithm}{Algorithms}

\usepackage[labelfont=bf]{caption}
\usepackage{svg}
